% How Many Backtest Winners Survive Deflation?
% Controlled study of DSR + multiple-testing haircuts.
% Every number in this manuscript is verified against results/results.json by
% scripts/check_paper_numbers.py; do not edit numbers without re-running it.
\documentclass[11pt]{article}
\usepackage[margin=1in]{geometry}
\usepackage{amsmath,amssymb}
\usepackage{booktabs}
\usepackage[round]{natbib}
\usepackage{microtype}
\usepackage[hidelinks]{hyperref}
\setlength{\emergencystretch}{2em}

\newcommand{\SRhat}{\widehat{SR}}
\newcommand{\E}{\mathbb{E}}
\newcommand{\V}{\mathbb{V}}
\newcommand{\Neff}{N_{\mathrm{eff}}}

\title{How Many Backtest Winners Survive Deflation?\\
\large A Controlled Study of the Deflated Sharpe Ratio and Multiple-Testing Haircuts}
\author{Eugen Soloviov}
\date{July 2026}

\begin{document}
\maketitle

\begin{abstract}
A parameter search over a backtest grid is a multiple-testing machine: the best
of many zero-skill configurations reliably produces an impressive Sharpe ratio.
We run four seeded, fully reproducible experiments with known ground truth to
measure how well the standard corrections repair this. Under a global null
(selecting the best of 1000 independent zero-skill strategies), the naive
single-test verdict certifies a discovery in every repetition (false-discovery
rate 1.000) at an average best annualized Sharpe of 1.63, while the Deflated
Sharpe Ratio (DSR), the Harvey--Liu haircuts, and White's Reality Check bring
the false-discovery rate to 0.001--0.057, at or near the nominal
$\alpha = 0.05$. The deflated benchmark $SR_0$ averages 1.63 annualized ---
exactly the noise ceiling: the bar a strategy must clear is not zero but the
best Sharpe that luck alone buys for the given search size and sample length.
A planted-signal sweep shows DSR power rising from 0.005 to 1.000, with the
50\%-power point just above that ceiling (about 1.73 annualized). Two
realistic moving-average searches (640 correlated trials each, bootstrap tests
on 5000 resamples) then probe the correlated-search regime. On pure noise
every estimator rejects the seductive winner (annualized Sharpe 0.81). On a
genuine regime-switching edge, whose in-sample \emph{selected} maximum is an
annualized Sharpe of 3.92, the Reality Check ($p = 0.0024$), the studentized
SPA-type test ($p = 0.0038$), and the Harvey--Liu haircuts (0.148--0.198) all
certify the edge --- but DSR fed the \emph{raw} trial count wrongly rejects it
($0.748 < 0.95$). The natural fix, an effective number of trials, turns out
not to be a single number: five standard estimators computed on the same trial
matrix disagree by two orders of magnitude (1.6 to 370.0), and at the smallest
estimates the deflation is nearly inert --- there the correct verdicts are
inherited from the winners' un-deflated significance, not produced by
deflation. The honest deliverable is a robustness band: the real edge is
retained for every effective count below 144.8 (four of the five estimators,
including genuinely deflating mid-range bars of annualized 1.85--2.00), while
the noise winner is rejected at every effective count. Bootstrap-based tests
sidestep the choice entirely by resampling the whole search jointly.
\end{abstract}

\section{Introduction}
\label{sec:intro}

Every systematic trading study contains a hidden hypothesis-testing problem.
A researcher who sweeps a parameter grid, keeps the configuration with the
highest backtest Sharpe ratio, and then reports that Sharpe ratio as if it
were a single experiment has run a multiple-testing procedure while quoting
single-test statistics. The reported performance is the maximum of many noisy
estimates, and the maximum of many zero-mean draws is emphatically not
zero-mean. \citet{bailey2014pseudo} call the resulting literature
pseudo-mathematical; \citet{harveyliuzhu2016} document the same selection
mechanism operating at the scale of an entire academic field, with hundreds of
published factors mined from broadly the same data.

Several corrections are in wide circulation. \citet{baileylopezdeprado2012psr}
introduced the Probabilistic Sharpe Ratio (PSR), which converts a sample
Sharpe ratio into the probability that the true Sharpe ratio exceeds a chosen
benchmark, accounting for sample length, skewness, and kurtosis.
\citet{baileylopezdeprado2014dsr} extended it to the Deflated Sharpe Ratio
(DSR), which sets that benchmark to the expected maximum Sharpe ratio a search
of $N$ independent zero-skill trials would produce --- so the winner is
measured against the best that luck alone could have delivered.
\citet{harveyliu2015backtesting} approach the same problem through classical
multiple-testing adjustments (Bonferroni, Holm, Benjamini--Hochberg--Yekutieli),
translating adjusted $p$-values back into a ``haircut'' on the reported Sharpe
ratio, with the economic framing developed in
\citet{harveyliu2014evaluating}. A third lineage tests the entire search at
once: White's Reality Check \citep{white2000reality} bootstraps the
distribution of the \emph{best} performance across all trials under the null,
and \citet{hansen2005spa} sharpened it against irrelevant alternatives via
studentization. \citet{lopezdeprado2018advances} surveys these tools alongside
the non-parametric Probability of Backtest Overfitting of
\citet{baileyetal2017pbo}.

What is harder to find in this literature is a \emph{controlled} comparison:
experiments in which the ground truth is known by construction --- because the
data-generating process is synthetic and seeded --- so that false-discovery
rates and detection power can be measured directly rather than argued about.
That is the gap this paper fills. We make no claim of new estimators; the
contribution is calibration evidence and an honest account of a failure mode
and of what can, and cannot, be salvaged from it.

Concretely, we run four experiments (Section~\ref{sec:design}) against the
estimators of Section~\ref{sec:estimators}:

\begin{enumerate}
\item \textbf{Null calibration.} Select the best of $N = 1000$ independent
  zero-skill strategies, 2000 times. The naive single test certifies a
  discovery in every repetition; DSR, the Harvey--Liu haircuts, and the
  Reality Check bring the false-discovery rate to between 0.001 and 0.057, at
  or near the nominal level. The deflated benchmark $SR_0$ reproduces the
  empirical noise ceiling to two decimals (both 1.63 annualized).
\item \textbf{Planted power.} With 25 genuine strategies hidden among 1000,
  DSR's detection power traces an S-curve that crosses just above the noise
  ceiling: edges below the ceiling are statistically indistinguishable from
  luck; edges above it are retained with power approaching one, at a measured
  false-positive rate of 0.000.
\item \textbf{A realistic search on noise.} A moving-average crossover grid of
  640 correlated configurations, applied to a pure random walk, yields a
  winner with annualized Sharpe 0.81. Every estimator rejects it, at every
  choice of effective trial count.
\item \textbf{The same search on a real edge.} On a regime-switching series
  the in-sample selected maximum is an annualized Sharpe of 3.92 and the
  bootstrap tests certify it decisively --- but DSR, fed the raw trial count
  of 640, \emph{wrongly rejects} the genuine edge. The obvious repair, an
  effective number of trials, is not a single number: five standard
  estimators computed on the same trial matrix span 1.6 to 370.0. Four of the
  five retain the edge --- including mid-range estimators whose deflated bar
  is a genuine annualized 1.85--2.00 --- and DSR retains it for every
  effective count below 144.8 while rejecting the noise winner at every
  effective count. We argue that this robustness band, not any point
  estimate, is the honest deliverable of effective-$N$ deflation, and we
  document where the band's lower anchor is misleading.
\end{enumerate}

Section~\ref{sec:results} reports the numbers, Section~\ref{sec:discussion}
draws the practical lessons, and Section~\ref{sec:limitations} states the
limitations plainly. All results derive from one seeded run of one public
harness; a companion script verifies every numeric claim in this manuscript
against the run's JSON output.

\section{Estimators}
\label{sec:estimators}

\paragraph{Notation.} Sharpe ratios are per-observation, $\SRhat = \hat\mu /
\hat\sigma$ computed on $T$ return observations, unless explicitly annualized;
annualized values multiply the per-observation value by $\sqrt{252}$ (we treat
252 observations as one year of daily data). $\Phi$ and $\Phi^{-1}$ denote the
standard normal CDF and quantile function. Tests are one-sided (skill means
$SR > 0$) at level $\alpha = 0.05$ throughout. Equation-number citations below
refer to the numbering printed in the primary sources, which we re-extracted
from the authors' posted full texts for
\citet{baileylopezdeprado2012psr,baileylopezdeprado2014dsr,harveyliu2015backtesting};
the tests of Section~\ref{sec:rcspa} are described per their standard
characterization, as discussed there.

\subsection{Probabilistic Sharpe Ratio}
\label{sec:psr}

The PSR of \citet{baileylopezdeprado2012psr}, their Eq.~(11), with the
estimated variance of $\SRhat$ from their Eq.~(8) (a result they attribute to
Mertens), is the probability that the true Sharpe ratio exceeds a benchmark
$SR^{*}$ given the sample evidence:
\begin{equation}
\label{eq:psr}
\widehat{\mathrm{PSR}}(SR^{*})
= \Phi\!\left[
\frac{(\SRhat - SR^{*})\,\sqrt{T-1}}
{\sqrt{\,1-\hat{\gamma}_{3}\SRhat+\frac{\hat{\gamma}_{4}-1}{4}\SRhat^{2}\,}}
\right],
\end{equation}
where $\hat\gamma_3$ is the sample skewness and $\hat\gamma_4$ the sample
kurtosis of the strategy's returns in the \emph{non-excess} convention
(normal returns give $\hat\gamma_4 = 3$). The convention matters: for normal
returns the denominator reduces to the classical asymptotic
$\V[\SRhat]\approx(1+SR^{2}/2)/T$, which only happens with non-excess
kurtosis; substituting excess kurtosis silently corrupts the deflation. All
quantities are computed in the returns' native frequency --- the source is
explicit that PSR is invariant to calendar conventions, so no annualization
enters the test itself. The same expression is restated as Eq.~(2) of
\citet{baileylopezdeprado2014dsr}.

\subsection{Deflated Sharpe Ratio}
\label{sec:dsr}

The DSR of \citet{baileylopezdeprado2014dsr} is PSR evaluated at a benchmark
that is not zero but the expected maximum Sharpe ratio that a search of $N$
independent zero-skill trials would produce. Let $\{\SRhat_{n}\}_{n=1}^{N}$ be
the Sharpe estimates of all trials in the search, with empirical mean
$\hat\mu_{SR}$ and variance $\widehat{V}_{SR}$. Their Eq.~(1) (derived in
their appendix as Eqs.~(5)--(6) via an extreme-value approximation to the
expected maximum of $N$ iid standard normal draws) gives the deflated
benchmark
\begin{equation}
\label{eq:sr0}
SR_{0} = \hat\mu_{SR} + \sqrt{\widehat{V}_{SR}}
\left[(1-\gamma)\,\Phi^{-1}\!\Big(1-\frac{1}{N}\Big)
+ \gamma\,\Phi^{-1}\!\Big(1-\frac{1}{Ne}\Big)\right],
\end{equation}
where $\gamma \approx 0.5772$ is the Euler--Mascheroni constant and $e$ is
Euler's number. Their Eq.~(1) states the benchmark under the null
$\E[\SRhat_n] = 0$, in which the mean term vanishes; we carry the empirical
mean $\hat\mu_{SR}$ to match the authors' own reference implementation
(function \texttt{getExpMaxSR} in the published paper), which does the same.
Then
\begin{equation}
\label{eq:dsr}
\widehat{\mathrm{DSR}} = \widehat{\mathrm{PSR}}(SR_{0}),
\end{equation}
their Eq.~(2), computed with the \emph{selected} strategy's own $\SRhat$, $T$,
$\hat\gamma_3$, and $\hat\gamma_4$ in Eq.~\eqref{eq:psr}. We flag a discovery
when $\widehat{\mathrm{DSR}} > 1-\alpha = 0.95$.

Three properties of Eq.~\eqref{eq:sr0} drive everything that follows. First,
$SR_0$ grows with the search: roughly like the square root of $\log N$ through
the quantile terms, and proportionally to the cross-trial dispersion
$\sqrt{\widehat{V}_{SR}}$, which itself shrinks like one over the square root
of $T$ under the null. The deflated bar is therefore a \emph{noise ceiling}
specific to the search size and sample length. Second, the dispersion term
treats the observed spread of trial Sharpes as pure luck: when part of that
spread is produced by genuinely skilled trials, Eq.~\eqref{eq:sr0} reads real
skill as more luck to deflate against. Third --- the caveat printed in the
source itself --- the derivation assumes the $N$ trials are
\emph{independent}. The paper's Appendix~3 addresses determining $N$ when they
are not; Section~\ref{sec:neff} returns to this point, which turns out to be
the single most consequential implementation decision in our experiments.

\subsection{Harvey--Liu haircuts}
\label{sec:hl}

\citet{harveyliu2015backtesting} start from the exact algebraic link between
the Sharpe ratio and the single-test $t$-statistic (their Eqs.~(1)--(2)):
$t = \SRhat\sqrt{T}$, from which a one-sided $p$-value follows,
$p = 1 - \Phi(t)$ under the normal approximation. Testing $M$ strategies
yields ordered $p$-values $p_{(1)} \le \dots \le p_{(M)}$, which are adjusted
by one of three classical procedures --- Bonferroni, Holm \citep{holm1979},
and Benjamini--Hochberg--Yekutieli
\citep[BHY;][]{benjaminihochberg1995,benjaminiyekutieli2001}:
\begin{align}
\text{Bonferroni:}\quad
& p^{\mathrm{Bonf}}_{(i)} = \min\!\big[M\,p_{(i)},\,1\big],
\label{eq:bonf}\\[2pt]
\text{Holm:}\quad
& p^{\mathrm{Holm}}_{(i)} = \min\!\Big[\max_{j\le i}\,(M-j+1)\,p_{(j)},\,1\Big],
\label{eq:holm}\\[2pt]
\text{BHY:}\quad
& p^{\mathrm{BHY}}_{(M)} = p_{(M)}, \qquad
p^{\mathrm{BHY}}_{(i)} = \min\!\Big[p^{\mathrm{BHY}}_{(i+1)},\,
\frac{M\,c(M)}{i}\,p_{(i)}\Big],
\label{eq:bhy}
\end{align}
where the step-up recursion in Eq.~\eqref{eq:bhy} runs from $i = M-1$ down to
the smallest $p$-value, with the normalizing constant \emph{multiplying} the
adjustment,
\begin{equation}
\label{eq:cm}
c(M) = \sum_{j=1}^{M}\frac{1}{j}.
\end{equation}
\citet{benjaminihochberg1995} originally set $c(M) = 1$, valid under
independence or positive dependence; Harvey and Liu adopt the
\citet{benjaminiyekutieli2001} harmonic-sum choice precisely because it
controls the false-discovery rate under \emph{arbitrary} dependence among the
test statistics --- the situation of correlated backtest trials. The constant
grows slowly: $c(1000) \approx 7.49$, so the price of dependence-robustness is
less than one order of magnitude. Bonferroni and Holm control the familywise
error rate $\mathrm{FWER} = \Pr(\text{at least one false rejection})$; BHY
controls the false-discovery rate
$\mathrm{FDR} = \E[\text{false rejections}/\text{rejections}]$, a more lenient
criterion for the \emph{bulk} of hypotheses. One rank-specific subtlety
matters for everything we report: the best-strategy haircut evaluates only the
\emph{top-ranked} hypothesis, and at rank one Holm coincides with Bonferroni
exactly --- $(M-1+1)\,p_{(1)} = M\,p_{(1)}$ --- so the two never disagree
about a search winner and should not be read as independent corroboration.
Moreover, BHY multiplies the top $p$-value by $M\,c(M) \ge M$, so for the top
pick BHY is the \emph{most} conservative of the three, even though it
controls the more lenient criterion overall.

The adjusted $p$-value of the search winner is then converted back through the
$t$-statistic link into a haircut Sharpe ratio $SR_{\mathrm{adj}}$ --- ``the
Sharpe ratio that would have resulted from a single test'' --- and the
headline quantity is the haircut fraction
\begin{equation}
\label{eq:haircut}
\mathrm{hc} = 1 - \frac{SR_{\mathrm{adj}}}{\SRhat}.
\end{equation}

\subsection{Reality Check and the studentized (SPA-type) variant}
\label{sec:rcspa}

The tests above operate on summary statistics of the trials. White's Reality
Check \citep{white2000reality} instead tests the whole search at once, using
the full $T \times K$ matrix of trial returns (in excess of a benchmark; ours
is zero, i.e.\ not trading). With $\bar f_{k}$ the sample mean performance of
trial $k$, the null is $H_{0}\colon \max_{k}\E[f_{k}] \le 0$ and the statistic
is the scaled best average, $V = \max_{k}\sqrt{T}\,\bar f_{k}$. Its
distribution under the null is approximated by the stationary bootstrap of
\citet{politisromano1994}: resampled series are built from blocks of
geometrically distributed random length (expected block length $1/p$),
wrapping circularly so the resample is itself stationary. Each bootstrap
replication recenters every trial by its own full-sample mean, so the
resampled world satisfies the null by construction, and the $p$-value is the
add-one fraction of bootstrap statistics
$V^{*b} = \max_{k}\sqrt{T}\,(\bar f^{*b}_{k} - \bar f_{k})$ that reach the
observed $V$, i.e.\ $\hat p = (1 + \#\{V^{*b} \ge V\})/(B+1)$. Because the
bootstrap resamples all $K$ trial series \emph{jointly}, the cross-trial
correlation structure is preserved exactly; no independent-trials assumption
enters anywhere.

\citet{hansen2005spa} showed that the unstandardized maximum can be dominated
by poor, high-variance alternatives, making the Reality Check conservative,
and proposed the Superior Predictive Ability (SPA) test with two
modifications: studentizing each trial by an estimate of its own standard
error before taking the maximum, and a sample-dependent (consistent)
recentering of the null distribution. Our implementation applies only the
first modification --- each trial's statistic is divided by its sample
standard deviation, with White's full recentering retained --- so it is a
\emph{studentized Reality Check}, not Hansen's full SPA; we label it
``SPA-type'' everywhere in this paper, and note that omitting Hansen's
recentering can only make it more conservative. One honesty flag: unlike
Sections~\ref{sec:psr}--\ref{sec:hl}, whose equations we re-extracted from the
primary texts, this subsection follows the standard characterization of
\citet{white2000reality}, \citet{hansen2005spa}, and
\citet{politisromano1994}; we cite no equation numbers from those papers.

\subsection{Effective number of trials: a spread, not a number}
\label{sec:neff}

Equation~\eqref{eq:sr0} needs the number of \emph{independent} trials, and a
parameter grid does not supply that number directly: neighboring
configurations trade nearly identical positions. Three strands of the
literature converge on the same warning. \citet{baileylopezdeprado2014dsr}
state the independence assumption and point to their Appendix~3 for an
effective-$N$ construction under dependence. \citet{harveyliu2015backtesting}
adopt the BHY constant of Eq.~\eqref{eq:cm} specifically for dependence
robustness. \citet{harveyliuzhu2016} model the correlation among the factors
they audit rather than counting them, arriving at a $t$-hurdle of roughly
three for contemporary discoveries.

There is, however, no canonical estimator of the effective trial count, and
the candidates disagree --- in our experiments, by two orders of magnitude on
the same matrix. We therefore compute five standard ones and report the whole
spread. With $\bar\rho$ the average pairwise correlation among the $K$ trial
return series, the average-correlation estimate is
\begin{equation}
\label{eq:neff}
\Neff^{\mathrm{corr}} = \frac{K}{1 + (K-1)\,\bar\rho},
\end{equation}
the variance-reduction factor of the \emph{mean} of $K$ equicorrelated
variables --- exact in that setting, and the smallest of the five in
practice. A functional mismatch deserves flagging here: DSR's benchmark is an
expected \emph{maximum} (an extreme-value quantity), not a mean, and there is
no reason the effective count appropriate for averaging transfers to
extremes; Eq.~\eqref{eq:neff} is best read as a lower anchor. The remaining
four work from the eigenvalues $\lambda_1 \ge \lambda_2 \ge \dots$ of the
$K \times K$ trial correlation matrix (which sum to $K$): the
\emph{participation ratio} $(\sum_i\lambda_i)^2 / \sum_i\lambda_i^2$;
\emph{PCA-95}, the number of leading eigenvalues needed to explain 95\% of
total variance; the \emph{Kaiser criterion}, the count of eigenvalues
exceeding one; and the \emph{Cheverud--Nyholt} estimate
$1 + (K-1)\bigl(1 - \mathrm{Var}(\lambda)/K\bigr)$, a standard recipe
imported from the multiple-testing literature in statistical genetics, known
to over-count effective tests under strong equicorrelation and in practice
the largest of the five (the upper anchor). We cite no primary sources for
these four; they are used here as widely known recipes, not endorsed
estimators.

Because $\widehat{\mathrm{DSR}}$ is monotonically decreasing in $N$, feeding
each estimate into Eq.~\eqref{eq:sr0} traces a curve, and the entire analysis
can be compressed into one number found by bisection:
\begin{equation}
\label{eq:nstar}
N^{*} = \max\bigl\{N \ge 2 : \widehat{\mathrm{DSR}}(N) > 1-\alpha\bigr\},
\end{equation}
the largest effective trial count at which the winner still survives. A
claimed discovery then comes with a robustness band --- ``retained for every
$\Neff < N^{*}$'' --- instead of a verdict conditioned on one contestable
point estimate. One domain note: the expected-maximum formula is defined for
$N \ge 2$ (the $\Phi^{-1}(1-1/N)$ term diverges as $N \to 1$), so fractional
estimates below two are evaluated at the $N = 2$ floor.

\section{Experimental design}
\label{sec:design}

All experiments are generated and analyzed by a single Python harness
(\texttt{scripts/run\_all.py}, estimators in \texttt{scripts/deflate.py}),
under Python 3.14.6 with NumPy 2.4.3. Every random stream is seeded, so
every number in this paper is bit-reproducible by one command. Ground truth is
known by construction in each experiment: we know which strategies have skill
because we planted them. A verification script
(\texttt{scripts/check\_paper\_numbers.py}) checks every numeric claim in this
manuscript against the harness output \texttt{results/results.json}; values
quoted below are rounded from that file, and quantities labeled
\emph{derived} are arithmetic combinations of its entries (annualization by
$\sqrt{252}$, differences, products, interpolation).

\subsection{Experiment 1: null calibration}
\label{sec:design-null}

Each repetition draws a $T \times N$ matrix of iid standard normal returns
with $T = 1000$ observations (about four years of daily data) and $N = 1000$
strategies --- every true Sharpe ratio is exactly zero. We select the winner
$\SRhat_{\max} = \max_n \SRhat_n$ and ask each procedure whether it certifies
a discovery at $\alpha = 0.05$:
\begin{itemize}
\item \emph{Naive:} one-sided single test on the winner,
  $p = 1-\Phi(\SRhat_{\max}\sqrt{T}) < \alpha$, ignoring the search;
\item \emph{DSR:} Eq.~\eqref{eq:dsr} $> 0.95$, with all $N$ trial Sharpes
  informing $SR_0$;
\item \emph{Harvey--Liu:} adjusted $p$-value of the winner below $\alpha$
  under each of Eqs.~\eqref{eq:bonf}--\eqref{eq:bhy} with $M = N$;
\item \emph{Reality Check:} bootstrap $p < \alpha$ on the full return matrix.
\end{itemize}
The closed-form procedures run over $M = 2000$ repetitions. The Reality
Check, being bootstrap-based, runs over 400 repetitions with 500 stationary
bootstrap resamples each (expected block length 20; the null data are iid, so
blocking is harmless). Under a global null every certified discovery is
false, so the empirical discovery rate \emph{is} the false-discovery rate.

\subsection{Experiment 2: planted power}
\label{sec:design-power}

Calibration is cheap if a test never fires; the counterpart question is power.
We repeat the setup of Experiment~1 but add a constant per-observation drift
$s_{\mathrm{true}}$ to $n_{\mathrm{true}} = 25$ randomly chosen strategies per
repetition, giving them true per-observation Sharpe $s_{\mathrm{true}} \in
\{0.05, 0.08, 0.12, 0.16, 0.20\}$ (annualized 0.79 to 3.17), with 1000
repetitions per level. We record whether DSR certifies the search winner, and
whether that winner is genuinely one of the planted strategies (a true
positive) or an impostor (a false positive). For reference we also record how
often the naive test certifies a winner that happens to be planted. Two scope
notes, stated up front: the power figures are specific to this configuration
(25 planted among $N = 1000$, $T = 1000$) --- we did not sweep these
parameters --- and Experiment~2 measures the power of DSR only; the retention
behavior of the other procedures is probed through the single case study of
Experiment~4.

\subsection{Experiments 3--4: a realistic search, without and with an edge}
\label{sec:design-search}

The first two experiments use independent trials --- exactly the world
Eq.~\eqref{eq:sr0} assumes. Real searches are not like that, so we build one.
A moving-average crossover strategy computes fast and slow simple moving
averages of the log price and holds the position
$\mathrm{sign}(\mathrm{MA}_{f} - \mathrm{MA}_{s}) \in \{-1, 0, +1\}$; the grid
sweeps $f \in \{3, 6, \ldots, 48\}$ (16 values) and
$s \in \{55, 60, \ldots, 250\}$ (40 values), $K = 640$ configurations in all.
Positions are strictly causal --- the position formed at observation $t$ earns
the return at $t+1$ --- and frictionless (no transaction costs; the point here
is selection inference, not net profitability). The price series has
$T = 756$ daily returns (three years at 252 observations per year), leaving
755 tradable strategy returns per configuration.

Two data-generating processes, identical in every respect except the presence
of skill to find:
\begin{itemize}
\item \emph{Noise (Experiment 3):} iid Gaussian returns with zero mean and
  standard deviation 0.01 per observation --- a pure random walk in log price.
  No configuration has true skill; the correct verdict is rejection.
\item \emph{Regime edge (Experiment 4):} the same Gaussian noise plus a
  regime-switching drift $d_t = s_t \cdot 0.006$, where the state
  $s_t \in \{-1, +1\}$ flips with probability 0.02 per observation (expected
  regime length 50 observations). Trends persist for weeks at a time; a
  crossover rule genuinely can extract this. An omniscient regime-follower
  would earn per-observation Sharpe 0.6 (about 9.5 annualized --- derived), so
  a good deal of real skill is available; the correct verdict is retention.
\end{itemize}
Each search is analyzed exactly as a practitioner would analyze a live sweep:
naive $p$-value and un-deflated PSR against zero for the winner; DSR under
the raw trial count $N = K$ and under each of the five effective-count
estimators of Section~\ref{sec:neff}, together with the survival crossing
$N^{*}$ of Eq.~\eqref{eq:nstar}; Reality Check and the SPA-type test on the
full $755 \times 640$ return matrix (5000 resamples for these case studies,
expected block length 20); and the three Harvey--Liu haircuts with $M = K$.

\section{Results}
\label{sec:results}

\subsection{Null calibration: the naive verdict is always wrong}
\label{sec:res-null}

Table~\ref{tab:null} shows the false-discovery rates. Selecting the best of
1000 zero-skill strategies produced an average best per-observation Sharpe
corresponding to an annualized 1.63, and the naive analyst's median $p$-value
on that winner was 0.00069 --- apparently overwhelming evidence, in every
single repetition (false-discovery rate 1.000). This is the multiple-testing
machine at work: the naive test answers ``could \emph{this one} strategy be
zero-skill?'' when the question posed by the search is ``could the \emph{best
of a thousand} be the product of zero skill?''

\begin{table}[t]
\centering
\caption{Experiment 1 (global null): empirical false-discovery rate at
$\alpha = 0.05$, i.e.\ the share of repetitions in which each procedure
certifies the best of $N = 1000$ zero-skill strategies as a discovery.
Closed-form procedures: 2000 repetitions; Reality Check: 400 repetitions of
500 stationary-bootstrap resamples. Bonferroni and Holm are reported as one
row because they coincide exactly for the top-ranked winner
(Section~\ref{sec:hl}). Monte Carlo standard error of a proportion near the
nominal level is about 0.005.}
\label{tab:null}
\begin{tabular}{llll}
\toprule
Procedure & Discovery rule & Error criterion & FDR \\
\midrule
Naive single test & $p < 0.05$ & none (ignores search) & 1.000 \\
DSR & $\widehat{\mathrm{DSR}} > 0.95$ & deflated benchmark & 0.001 \\
Harvey--Liu, Bonferroni $=$ Holm & adj.\ $p < 0.05$ & FWER & 0.057 \\
Harvey--Liu, BHY & adj.\ $p < 0.05$ & FDR & 0.007 \\
White Reality Check & bootstrap $p < 0.05$ & test of the max & 0.0225 \\
\bottomrule
\end{tabular}
\end{table}

Every correction restores control, each with its own personality. Bonferroni
and Holm --- identical for the winner by construction --- sit at 0.057,
statistically indistinguishable from the nominal level (within about two
Monte Carlo standard errors of 0.005; the small excess is consistent with the
normal approximation to the $t$-statistic being slightly optimistic in the
extreme tail where $\min_n p_n$ lives). BHY lands at 0.007 --- close to
$\alpha / c(1000) \approx 0.007$ (derived), which is exactly the conservatism
its harmonic-sum constant buys as insurance against dependence that this iid
experiment does not actually contain, and consistent with BHY being the most
conservative of the three at rank one. The Reality Check comes in at 0.0225,
conservative in the direction \citet{hansen2005spa} predicted for the
unstandardized statistic. DSR is the strictest of all at 0.001, and the
reason is instructive: its bar is not a significance level but the noise
ceiling itself. The average deflated benchmark $SR_0$ across repetitions was
1.63 annualized (derived from the per-observation value) --- identical to two
decimals with the average realized maximum. Eq.~\eqref{eq:sr0} is doing
precisely its job: predicting what the best of $N$ lucky draws looks like,
and refusing to be impressed by it. Consistent with calibration, the DSR
value of the null winner averaged 0.495 across repetitions, indistinguishable
from the ideal 0.5 --- the estimator judges the typical lucky winner exactly
as likely as not to beat the luck benchmark, which is the correct answer.

\subsection{Planted power: the S-curve crosses just above the noise ceiling}
\label{sec:res-power}

Table~\ref{tab:power} reports DSR's detection power against the planted true
Sharpe ratio. The pattern is a sharp S-curve positioned just above the noise
ceiling of Experiment~1: linear interpolation between the two adjacent grid
levels puts the 50\%-power point near an annualized 1.73 (derived), slightly
above the ceiling of 1.63 rather than at it. A genuine annualized Sharpe of
0.79 --- respectable in production --- is essentially undetectable (power
0.005): it lives so far below the ceiling that the search's own noise drowns
it, and indeed at that level the search winner is one of the planted
strategies in only about two-thirds of repetitions (0.670). At annualized
1.27 the power is still only 0.090. Crossing the ceiling changes everything:
0.651 at annualized 1.90, then 0.998 at 2.54 and 1.000 at 3.17. Throughout
the sweep the measured false-positive rate of DSR is 0.000 --- in no
repetition at any level did DSR certify an impostor. The naive column is a
foil: it fires on essentially every winner (compare Table~\ref{tab:null}), so
its apparent sensitivity carries no information about skill. We reiterate the
scope note from Section~\ref{sec:design-power}: these are DSR power figures
at one configuration, not a general power study.

\begin{table}[t]
\centering
\caption{Experiment 2 (planted signal): DSR detection power over 1000
repetitions per level, with 25 planted strategies among $N = 1000$ and
$T = 1000$. ``Naive hit'' is the share of repetitions in which the naive test
certifies the winner \emph{and} the winner is genuinely planted; since the
naive test essentially always fires, this column mostly measures how often
the best-by-Sharpe trial is a planted one.}
\label{tab:power}
\begin{tabular}{rrrrr}
\toprule
True $SR$ (per obs.) & True $SR$ (annual) & DSR power & DSR false-pos.\ rate & Naive hit \\
\midrule
0.05 & 0.79 & 0.005 & 0.000 & 0.670 \\
0.08 & 1.27 & 0.090 & 0.000 & 0.984 \\
0.12 & 1.90 & 0.651 & 0.000 & 1.000 \\
0.16 & 2.54 & 0.998 & 0.000 & 1.000 \\
0.20 & 3.17 & 1.000 & 0.000 & 1.000 \\
\bottomrule
\end{tabular}
\end{table}

The practical reading is sobering and useful: given this search size and
sample length, the deflated bar sits at an annualized Sharpe of roughly 1.63,
and \emph{no statistical procedure operating on the same evidence can do
fundamentally better} --- a sub-ceiling edge is not distinguishable from the
best of a thousand lucky draws, however real it is. Deflation does not
destroy value; it prices the evidence.

\subsection{A realistic search on noise: everything correctly rejects}
\label{sec:res-noise}

Table~\ref{tab:searches} (left column) reports the moving-average sweep on a
pure random walk. The winner --- crossover $(45, 120)$ --- earned an
annualized Sharpe of 0.81 over three years. Plotted, such an equity curve
looks eminently fundable. Every layer of scrutiny disagrees. Even before any
deflation, the winner is not single-test significant: the naive $p$-value is
0.081 and the un-deflated PSR against a zero benchmark is 0.918, short of the
0.95 threshold (the sample is short and the Sharpe modest). Deflation-aware
diagnostics reject with progressively more context. DSR with the raw
$N = 640$ gives 0.431. The Reality Check and the SPA-type test, which
resample the entire 640-trial matrix jointly (5000 resamples), return
$p$-values of 0.570 and 0.569: the observed best is entirely typical of a
searched random walk. The Harvey--Liu adjusted $p$-value of the winner is
1.00 under Holm, and all three haircuts are total (1.00): the certified
Sharpe after multiple-testing adjustment is zero.

The effective-count band (Table~\ref{tab:neff}, left half) adds the
robustness statement: the five estimators span 1.6 to 379.9, the implied
deflated bars run from an annualized 0.31 to 0.87, and the implied DSR values
run from 0.805 down to 0.455 --- \emph{every} choice rejects, and the
survival crossing of Eq.~\eqref{eq:nstar} confirms there is no effective
trial count at which this winner survives. One honesty note that
Section~\ref{sec:discussion} develops: at the smallest estimates the bar has
collapsed to near zero and DSR nearly coincides with the un-deflated PSR
(0.918), so the rejection \emph{there} rests on the winner's own modest raw
significance rather than on any real deflation; the genuinely deflating
verdicts are the mid-range and upper rows.

\begin{table}[t]
\centering
\caption{Experiments 3--4: the same 640-configuration moving-average
crossover search ($T = 755$ tradable returns) on a pure random walk (left)
and on a regime-switching series with a real trend edge (right). Winner
Sharpe ratios are in-sample \emph{selected} maxima and therefore
selection-inflated by construction. Bootstrap tests use 5000 resamples.
Bonferroni and Holm haircuts coincide for the top-ranked winner
(Section~\ref{sec:hl}) and are reported as one row. The DSR discovery
threshold is 0.95; for RC, SPA-type, and adjusted $p$-values the threshold is
$\alpha = 0.05$.}
\label{tab:searches}
\begin{tabular}{lrr}
\toprule
 & Noise (no skill) & Regime edge (real skill) \\
\midrule
Winner $(f, s)$ & $(45, 120)$ & $(3, 55)$ \\
Winner $SR$, in-sample selected (annualized) & 0.81 & 3.92 \\
Naive one-sided $p$ & 0.081 & $6.1\times 10^{-12}$ \\
Un-deflated $\widehat{\mathrm{PSR}}(0)$ & 0.918 & 1.000 \\
Mean pairwise correlation $\bar\rho$ & 0.61 & 0.62 \\
\addlinespace
$SR_0$ with raw $N = 640$ (annualized) & 0.91 & 3.51 \\
DSR with raw $N$ & 0.431 & 0.748 \\
\addlinespace
Reality Check $p$ & 0.570 & 0.0024 \\
SPA-type (studentized RC) $p$ & 0.569 & 0.0038 \\
\addlinespace
Harvey--Liu haircut, Bonferroni $=$ Holm & 1.00 & 0.148 \\
Harvey--Liu haircut, BHY & 1.00 & 0.198 \\
Holm-adjusted $p$ of winner & 1.00 & $3.9\times 10^{-9}$ \\
\addlinespace
Verdict, DSR (raw $N$) & reject (correct) & reject (\textbf{wrong}) \\
Verdict, RC and SPA-type & reject (correct) & retain (correct) \\
Verdict, Harvey--Liu (all three) & reject (correct) & retain (correct) \\
\bottomrule
\end{tabular}
\end{table}

\begin{table}[t]
\centering
\caption{The effective-trial-count band for both searches: five standard
estimators of $\Neff$, the deflated benchmark and DSR each implies, and the
survival crossing $N^{*}$ of Eq.~\eqref{eq:nstar}. Annualized $SR_0$ values
are reported by the harness; the raw grid count is included for reference.
The average-correlation estimates fall below the $N = 2$ domain floor of
Eq.~\eqref{eq:sr0} and are evaluated at that floor. The five estimators span
two orders of magnitude on the same matrix --- the reason we report a band
rather than a point verdict.}
\label{tab:neff}
\small
\begin{tabular}{lrrrcrrr}
\toprule
 & \multicolumn{3}{c}{Noise (no skill)} & & \multicolumn{3}{c}{Regime edge (real skill)} \\
\cmidrule{2-4}\cmidrule{6-8}
Estimator of $\Neff$ & $\Neff$ & $SR_0$ (ann.) & DSR & & $\Neff$ & $SR_0$ (ann.) & DSR \\
\midrule
Average correlation, Eq.~\eqref{eq:neff} & 1.6 & 0.31 & 0.805 & & 1.6 & 0.25 & 1.000 \\
Participation ratio & 2.4 & 0.35 & 0.785 & & 2.4 & 0.43 & 1.000 \\
PCA (95\% of variance) & 17 & 0.61 & 0.633 & & 16 & 1.85 & 1.000 \\
Kaiser criterion & 22 & 0.64 & 0.616 & & 21 & 2.00 & 0.999 \\
Cheverud--Nyholt & 379.9 & 0.87 & 0.455 & & 370.0 & 3.31 & 0.845 \\
Raw grid count $K$ & 640 & 0.91 & 0.431 & & 640 & 3.51 & 0.748 \\
\midrule
Survives (DSR $> 0.95$) for $\Neff < N^{*}$ &
\multicolumn{3}{c}{no $\Neff$ (never survives)} & &
\multicolumn{3}{c}{$N^{*} = 144.8$} \\
\bottomrule
\end{tabular}
\end{table}

\subsection{The same search on a real edge: the effective-$N$ band}
\label{sec:res-signal}

The right column of Table~\ref{tab:searches} is the reason this paper exists.
The regime-switching series contains genuine, exploitable trend persistence,
and the search finds it: the winner $(3, 55)$ has an in-sample selected
maximum of 3.92 annualized --- selection-inflated by construction, but
resting on real skill --- with a naive $p$-value of $6.1\times 10^{-12}$ and
an un-deflated PSR of 1.000. The bootstrap tests agree emphatically: with
5000 resamples the Reality Check returns $p = 0.0024$ and the SPA-type test
$p = 0.0038$ --- interior $p$-values, not the resolution floor of the
bootstrap. The Harvey--Liu haircuts are mild: 0.148 under Bonferroni and Holm
(one number, by the rank-one identity), 0.198 under BHY --- the most
conservative for a top pick, exactly as Section~\ref{sec:hl} predicts ---
leaving certified annualized Sharpes of 3.33 and 3.14 (derived), with a
Holm-adjusted $p$-value of $3.9\times 10^{-9}$. Real skill, correctly
retained, even after paying full freight for 640 trials.

DSR with the raw trial count reaches the opposite verdict. Fed $N = 640$,
Eq.~\eqref{eq:sr0} sets the deflated benchmark at an annualized 3.51 ---
demanding that the winner beat the expected best of 640 \emph{independent}
zero-skill trials --- and Eq.~\eqref{eq:dsr} returns 0.748, below the 0.95
threshold. A practitioner following the raw-$N$ recipe would discard a
genuine edge worth an annualized Sharpe of nearly four. Two distinct
mechanisms inflate that bar. The first is multiplicity: the grid's 640
trials have mean pairwise correlation 0.62 --- neighboring crossovers hold
nearly the same positions --- so 640 vastly overstates the number of
independent chances luck was given. The second is subtler: the dispersion
term of Eq.~\eqref{eq:sr0} treats the cross-trial spread of Sharpe estimates
as luck, but in this search much of the spread is real --- the per-observation
trial-Sharpe dispersion $\sqrt{\widehat{V}_{SR}}$ is 0.079 here versus 0.014
on the noise search (derived) because genuinely skilled regions of the grid
pull the distribution apart, and the formula reads that skill dispersion as
more luck to deflate against.

The effective-count band (Table~\ref{tab:neff}, right half) is the honest
repair, and it must be read whole. The five estimators again span two orders
of magnitude (1.6 to 370.0). Four of the five retain the edge. At the bottom
anchor (average correlation 1.6, evaluated at the $N = 2$ floor) the bar
collapses to an annualized 0.25 and DSR returns 1.000 --- but at that anchor
the deflation is nearly inert, and the retention is simply the winner's
un-deflated PSR (1.000) shining through. The informative rows are the
mid-range estimators: PCA-95 ($\Neff = 16$) and Kaiser ($\Neff = 21$) set
genuinely deflating bars of an annualized 1.85 and 2.00 --- comparable to the
noise ceiling of Experiment~1 --- and the edge still clears them decisively
(DSR 1.000 and 0.999). Only the near-independence Cheverud--Nyholt estimate
(370.0, close to the raw count of 640) flips the verdict, rejecting at 0.845.
The survival crossing summarizes the whole curve: this winner is retained for
every effective trial count below $N^{*} = 144.8$ --- more than a fifth of
the raw grid --- and the noise winner of Section~\ref{sec:res-noise} is
retained at none. That pair of statements, robust across every defensible
mid-range estimate, is what this experiment licenses; a single ``corrected''
DSR at one hand-picked $\Neff$ is not.

\section{Discussion}
\label{sec:discussion}

\subsection{Deflation restores calibration, and the bar is the noise ceiling}

The headline result of Experiments 1--2 is that the machinery works exactly
as advertised when its assumptions hold. The naive verdict on a searched
winner is not slightly optimistic --- it is always wrong under the null
(false-discovery rate 1.000), with median apparent significance of 0.00069.
Every principled correction restores the error rate to the neighborhood of
the nominal $\alpha = 0.05$ (0.001--0.057). And the deflated benchmark is not
an abstract penalty: $SR_0$ reproduced the empirically realized noise ceiling
to two decimals (1.63 annualized for a search of 1000 trials over 1000
observations). The practical form of this statement deserves emphasis in any
research process: \emph{a search has a Sharpe budget that luck will pay
regardless of skill}, computable in closed form from the trial count, trial
dispersion, and sample length, and no winner below that budget is evidence
of anything. Power against real edges is then a question of where the edge
sits relative to the ceiling (Table~\ref{tab:power}): negligible below it,
near-perfect above it, with the 50\%-power point just above the ceiling
(about an annualized 1.73, derived) and false positives at 0.000 throughout.

\subsection{Effective $N$: report a band, and anchor it honestly}
\label{sec:disc-neff}

Experiment 4 documents the failure mode that we believe practitioners most
need to know about, and Experiments 3--4 together show what can honestly be
salvaged. It is tempting --- it feels \emph{conservative} --- to feed DSR the
raw size of the parameter grid. On a correlated search this over-deflates
twice over: the raw count overstates the independent chances luck was given,
and the dispersion term simultaneously reads real-skill spread as luck
(0.079 versus 0.014 per observation across our two searches, derived). In our
regime-edge search the resulting benchmark (annualized 3.51) exceeded what
even a strong genuine edge could clear on three years of data, and DSR
rejected a strategy whose reality is not in doubt --- we planted it.

The obvious repair --- substitute an effective trial count --- immediately
runs into the fact that there is no such number, only estimators that
disagree by two orders of magnitude on the same matrix (1.6 to 379.9 on
noise, 1.6 to 370.0 on the edge). Worse, the band's ends are individually
treacherous. At the bottom, the average-correlation estimate
(Eq.~\eqref{eq:neff}) lands below two on both searches, where the
expected-maximum benchmark barely deflates at all: the bars are an annualized
0.25 and 0.31, and DSR effectively degenerates to the un-deflated PSR against
zero. Both verdicts at that anchor are \emph{inherited} --- the edge is
retained because its raw PSR is 1.000, and the noise winner is rejected
because its raw PSR happens to be 0.918, short of 0.95. The latter is luck of
the seed: this particular noise winner is modest (annualized 0.81). A
luckier noise draw, strong enough to clear single-test significance --- and
Experiment~1 shows searched noise winners routinely reach an annualized 1.63
--- would sail over a bar of 0.31 unchallenged. An effective count below two
does not deflate; it abdicates. We also note the functional mismatch of
Section~\ref{sec:neff}: Eq.~\eqref{eq:neff} is the variance-reduction factor
for a \emph{mean} of equicorrelated variables, imported into a formula about
the \emph{maximum}; nothing guarantees those two notions of ``effectively
independent'' coincide. At the top, the Cheverud--Nyholt estimate --- a
published, widely used recipe, known to over-count under strong
equicorrelation --- lands at 370.0 and reproduces the raw-count error almost
exactly (DSR 0.845 versus 0.748), rejecting the genuine edge. So the claim
``any defensible dependence adjustment beats the raw count'' is false, and we
do not make it.

What survives scrutiny is the band itself. The mid-range estimators ---
participation ratio, PCA-95, Kaiser, spanning roughly 2.4 to 22 --- impose
genuine deflation (bars up to an annualized 2.00 on the edge search) and
deliver the correct verdict on both searches. The survival crossing
$N^{*}$ compresses the whole curve into one auditable statement: the edge is
retained for every $\Neff < 144.8$, the noise winner for none. A discovery
claim of the form ``retained across the entire defensible range of effective
counts, with a survival margin documented'' is weaker than a single certified
number, but it is the strongest statement this evidence actually supports ---
and unlike the point estimate, it does not silently hinge on which
effective-count recipe the analyst happened to choose.

\subsection{What the bootstrap tests know that DSR does not}

The Reality Check and the SPA-type test never faced any of this because they
never needed a trial count: resampling the full return matrix jointly carries
the entire cross-trial dependence structure into the null distribution
automatically. Their separation on the two searches is stark --- $p = 0.570$
and 0.569 on noise versus $p = 0.0024$ and 0.0038 on the real edge, with 5000
resamples --- and required no judgment calls. This robustness has a price and
a scope. The price is data and computation: RC and SPA-type need the full
$T \times K$ matrix of trial returns and thousands of bootstrap replications,
whereas DSR needs only the list of trial Sharpe ratios --- often the only
thing archived from a historical search, and cheap at any scale. The scope
difference is subtler: the two families answer different questions. DSR asks,
``is the winner exceptional \emph{relative to this search's own luck
distribution}?'' RC and SPA-type ask, ``after accounting for the full search,
does the best rule beat the benchmark at all?'' The first is a
selection-adjusted quality bar on the champion; the second is an existence
test for skill anywhere in the family. On our clean constructions they agree
--- on noise everything rejects at every effective count, and on the real
edge the bootstrap tests, the Harvey--Liu haircuts, and DSR across its entire
defensible mid-range all retain (Table~\ref{tab:searches},
Table~\ref{tab:neff}) --- and their agreement across such different
mechanisms is itself evidence; in messier settings they are complements, not
substitutes.

\subsection{A practical recipe}

Our results suggest a concrete checklist for anyone auditing a backtest
search. Log every trial, not just the winner --- every estimator here
consumes the whole search, and the common practice of archiving only the
champion destroys exactly the information deflation needs. Report the
deflated benchmark $SR_0$ alongside the winner's Sharpe, so readers see the
noise ceiling the search itself created, and label the winner's Sharpe as the
in-sample selected maximum that it is. Do not report DSR at a single
effective trial count: compute the spread of standard estimators and the
survival crossing $N^{*}$, and distrust both ends of the band --- the
smallest estimates barely deflate (the verdict is just un-deflated PSR), and
the largest reproduce the raw-count over-deflation. A verdict that flips
inside the defensible mid-range is not a discovery. When trial return series
are available, run the Reality Check or its studentized SPA-type variant as
an assumption-light cross-check. And treat the Harvey--Liu haircut as the
communication layer: a statement like ``after Holm adjustment across 640
trials the certified Sharpe drops by 0.148'' is legible to any investment
committee in a way that a bootstrap $p$-value is not --- remembering that for
a single winner Holm \emph{is} Bonferroni, and BHY is the most conservative
of the three.

\section{Limitations}
\label{sec:limitations}

\begin{itemize}
\item \textbf{Synthetic data, by design.} Both data-generating processes ---
  iid Gaussian returns and a two-state regime-switching drift --- were chosen
  so that ground truth is known exactly, which is the point of a calibration
  study and also its boundary. Real returns have fat tails, volatility
  clustering, and structural breaks; our experiments say nothing about how
  large $T$ must be for Eq.~\eqref{eq:psr}'s moment corrections to absorb
  such features, only that the selection-bias machinery works when its
  assumptions hold.
\item \textbf{No consensus effective-$N$ estimator, and we do not supply
  one.} The five estimators we report are standard recipes, not derivations
  matched to DSR's extreme-value benchmark; the average-correlation form has
  the functional mismatch discussed in Section~\ref{sec:neff} (a
  mean-variance quantity applied to a maximum), and Cheverud--Nyholt is known
  to over-count under strong equicorrelation. Bailey and L\'opez de Prado's
  Appendix~3 clustering construction is the treatment we did not implement.
  Our claim is deliberately limited to the band: on these two searches the
  verdicts are stable across the defensible mid-range and summarized by
  $N^{*}$, and the band's ends are individually misleading in the ways
  documented in Section~\ref{sec:disc-neff}.
\item \textbf{Power evidence is DSR-only and configuration-specific.}
  Experiment~2 measures the power of DSR at one configuration
  ($n_{\mathrm{true}} = 25$, $N = 1000$, $T = 1000$) with no sensitivity
  sweep; the retention behavior of the Reality Check, the SPA-type test, and
  the haircuts is probed only through the single case study of Experiment~4.
\item \textbf{RC/SPA described, not re-derived; SPA-type is not Hansen's
  full test.} As flagged in Section~\ref{sec:rcspa}, we did not re-extract
  the equations of \citet{white2000reality}, \citet{hansen2005spa}, or
  \citet{politisromano1994} from the primary texts; our description follows
  their standard characterization. Our implementation is a studentized
  Reality Check: it adopts Hansen's studentization but retains White's full
  recentering, and does not implement Hansen's consistent (sample-dependent)
  recentering, so it is labeled SPA-type throughout and is, if anything,
  conservative relative to Hansen's SPA.
\item \textbf{Experiments 3--4 are case studies.} The two searches are single
  seeded draws --- deliberately, since they play the role of worked examples
  with known truth. The calibration and power claims rest on the
  Monte Carlo experiments (2000 and 1000 repetitions); the search
  experiments demonstrate existence of the raw-count failure and the shape of
  the effective-count band, not their frequency across seeds.
\item \textbf{No cost model, no execution claims.} The searches are
  frictionless and long--short unconstrained; wall-clock performance of the
  harness is likewise not a claim of this paper.
\end{itemize}

\section{Conclusion}
\label{sec:conclusion}

Under controlled conditions with known ground truth, the deflation toolkit
does what it promises. The naive single-test verdict on a searched winner is
worthless --- wrong with certainty under the null --- while the Deflated
Sharpe Ratio, the Harvey--Liu haircuts, and the Reality Check all restore
error control to the neighborhood of the nominal level, and the deflated
benchmark equals the noise ceiling of the search to two decimals. Power is a
function of where an edge sits relative to that ceiling, so deflation
converts an unanswerable question (``is this backtest good?'') into a
computable one (``does this edge clear what luck alone would have bought?'').

The toolkit's one sharp edge is the trial count inside DSR. On a realistic
correlated grid, the raw count rejected a genuine edge that every other
estimator certified --- and the seemingly obvious repair, an effective number
of trials, is not a number but a contested range spanning two orders of
magnitude, whose smallest values barely deflate at all and whose largest
reproduce the raw-count error. What the evidence does support is a
robustness band: on these searches the genuine edge survives every effective
count below 144.8, including mid-range bars that deflate as hard as the
noise ceiling itself, while the noise winner survives none. The bootstrap
tests sidestep the issue entirely by construction, at the cost of requiring
the full trial return matrix. Used together --- DSR reported as a band with
its survival crossing, RC and the SPA-type test for the search, a haircut for
the committee --- the estimators delivered a coherent verdict on every
question our experiments could pose, and their agreement across such
different mechanisms is the most reassuring calibration result of all.

\paragraph{Reproducibility.} All code, tests, and outputs accompany this
paper: \texttt{scripts/run\_all.py} regenerates
\texttt{results/results.json} from fixed seeds (Python 3.14.6, NumPy 2.4.3);
\texttt{scripts/\allowbreak check\_paper\_numbers.py} verifies every numeric
claim in this manuscript against that file and fails on any mismatch;
\texttt{tests/} contains deterministic invariant tests for every estimator
used.

\bibliographystyle{plainnat}
\bibliography{refs}

\end{document}
